\documentclass[10pt,a4paper]{article}
\usepackage[margin=2.2cm]{geometry}
\usepackage{fontspec}
\setmonofont{DejaVu Sans Mono}[Scale=0.84]
\newfontfamily\unifont{DejaVu Serif}[Scale=0.9]
\usepackage{amsmath,amssymb,mathtools}
\setcounter{MaxMatrixCols}{20}
\usepackage{longtable,booktabs,array}
\usepackage{fancyvrb}
\usepackage[dvipsnames]{xcolor}
\usepackage[colorlinks=true,linkcolor=NavyBlue,urlcolor=NavyBlue]{hyperref}
\usepackage{microtype}
\usepackage{needspace}
\usepackage{tocloft}\setlength{\cftbeforesecskip}{0.45em}
\setlength{\parindent}{0pt}\setlength{\parskip}{0.55em}
\setlength{\emergencystretch}{8em}
\makeatletter
\renewcommand\section{\@startsection{section}{1}{\z@}{-3.5ex \@plus -1ex \@minus -.2ex}{2.3ex \@plus.2ex}{\normalfont\Large\bfseries\raggedright}}
\renewcommand\subsection{\@startsection{subsection}{2}{\z@}{-3.25ex\@plus -1ex \@minus -.2ex}{1.5ex \@plus .2ex}{\normalfont\large\bfseries\raggedright}}
\renewcommand\subsubsection{\@startsection{subsubsection}{3}{\z@}{-3.25ex\@plus -1ex \@minus -.2ex}{1.5ex \@plus .2ex}{\normalfont\normalsize\bfseries\raggedright}}
\makeatother
\renewcommand{\arraystretch}{1.15}
\allowdisplaybreaks
\begin{document}
\title{dirac16complex as dark matter and dark energy: provenance\\[0.4em]\large Sources, data, decisions, derivations and the verification record of the analysis of section 24 of the notebook 7+1.nb}
\author{Repository: github.com/once-ere/7plus1 — prepared for Nash\\ Notebook: 7+1.nb (Wolfram Language 15.0.1)\\ Document generated by tools/make\_docs.py on 2026-10-28}
\date{}
\maketitle
\begin{abstract}
This record states what was asked, which results of the earlier sections and which observational data were used and where they come from, every decision taken with the equation behind it, how each result was derived and how it was checked: the 20 checks of section 24 of 7+1.nb (20 True), the 20 boolean checks of the independent Python program tools/ref\_dark.py (20 True), the comparison of the 29 groups of numbers printed by both programs, the program versions and the SHA-256 hashes of the files.
\end{abstract}
\tableofcontents
\Needspace{12\baselineskip}\section*{1. The task}\addcontentsline{toc}{section}{1. The task}
On 27 October 2026 the author asked for the following (summarized; the original wording is the author’s private working material and is not reproduced):

\begin{enumerate}
\item After computing accurate approximate solutions of the coupled field equations for the three gravity regimes (sections 20–23 of 7+1.nb), analyze the validity of associating dark matter (and dark energy) with dirac16complex, on the basis of those solutions.
\item Document the analysis in a documentation file with a companion provenance file and a symbol file, each in Markdown and LaTeX, the LaTeX compiled to PDF; never refer the reader to another file; put every equation that is created or used into the file that uses it.
\item Push every completed stage to github.com/once-ere/7plus1 (respecting the repository’s .gitignore) and verify the repository; continue through all stages without stopping.
\end{enumerate}
\textbf{Interpretation.} "Associating dark matter with dirac16complex" is read as the hypothesis that the cold dark matter of the universe consists of quanta of dirac16complex (its lightest Kaluza–Klein mode, or its Kaluza–Klein tower); "dark energy" as the hypothesis that dirac16complex — as matter, as a condensate or through its vacuum energy — causes the accelerated expansion. "Validity" is judged against (i) the solutions of the three regimes (what they derive exactly), (ii) order-of-magnitude estimates with natural couplings, and (iii) present data. Every statement is labelled DERIVED, ESTIMATE or OPEN, and the labels are part of the result.

The coordinates of 7+1.nb are kept: $X=(x_0,\dots,x_7)$, $x_0$ the time (the direction of the octonion unit $e_0=1$), $x_1,x_2,x_3$ ordinary space, $x_4,\dots,x_7$ the four auxiliary space-like coordinates.

\Needspace{8\baselineskip}\subsection*{1.1 The signature: the author’s explanation of 27 October 2026}\addcontentsline{toc}{subsection}{1.1 The signature: the author’s explanation of 27 October 2026}
\textbf{Source.} On 27 October 2026 the author of the project explained that the xAct list $\{7,1,0\}$ declares seven space-like dimensions and one time-like dimension, while $\mathrm{diag}(+1,-1,\dots,-1)$ is simply the squares of the octonion basis, and asked that this be made clear in all documentation, examples and provenance records. \textbf{What it corrected.} The version of 7+1.nb pushed in the commits 52d8729 to dc2f178 of github.com/once-ere/7plus1 declared the xAct metric of section 12 with \texttt{DefMetric[\{1, 7, 0\}, …]}, reading the list as the signs of $\eta$; since then section 12 declares \texttt{DefMetric[\{7, 1, 0\}, g71[-a, -b], CD, …]}, exactly as the author’s notebook Einstein-Lovelock-7+1-Nash.nb does. \textbf{What was added.} Checks in 7+1.nb — section 4.7: $\eta$ is the squares of the octonion basis, $e_Ae_A=\eta_{AA}e_0$ and $\mathrm{Re}(e_Ae_B)=\eta_{AB}$; section 12.2: \texttt{SignatureOfMetric[g71]} $=\{7,1,0\}$ with \texttt{SignDetOfMetric[g71]} $=-1$, $\mathrm{Re}(e_Ae_B)=\eta_{AB}$, and $\mathrm{Re}(x\,x)=\eta_{AB}x^Ax^B$; section 17: the flat metric $\mathrm{diag}(+1,-1,-1,-1)$ of the world brane is the squares of the basis $e_0,e_1,e_2,e_3$ of the quaternion subalgebra — and the same statements in the Python programs \texttt{tools/\allowbreak{}ref71.py}, \texttt{tools/\allowbreak{}ref\_\allowbreak{}gravity.py} and \texttt{tools/\allowbreak{}ref\_\allowbreak{}sm.py}; the explanation itself in the text of 7+1.nb (sections 1.2, 1.4, 4.2, 4.7, 7, 11, 12.2, 13, 15.1, 17) and in every document, symbol file and provenance file of the repository. \textbf{What did not change.} No result. Compared output by output with the backup below, every output of 7+1.nb is the same, except the declared list itself, the new checks and the check counts (and, as in every run, the name of the private context and the measured times); in particular every component of every Lovelock tensor is unchanged. The components follow from the component matrix of the metric, and the sign of the determinant, \texttt{SignDetOfMetric}, is $-1$ for both lists. \textbf{Backup.} The evaluated notebook as it was before is kept in the workspace as \texttt{backups/\allowbreak{}7+1\_\allowbreak{}2026-10-27\_\allowbreak{}before\_\allowbreak{}signature.nb} (the folder backups/ is excluded from the repository by .gitignore), 1034401 bytes, SHA-256 \texttt{06094248\allowbreak{}8caf4f45\allowbreak{}120ffd37\allowbreak{}be0123e9\allowbreak{}6992a31c\allowbreak{}d7fdbde0\allowbreak{}9ee4dc4c\allowbreak{}ed4f5d3a}; \texttt{tools/\allowbreak{}verify\_\allowbreak{}7p1\_\allowbreak{}nb.py} compares every input, text and section cell of 7+1.nb with it and finds every cell present, each one unchanged or changed by this explanation, and it compares every output in the same way.

\textbf{The explanation, in full.} The 7+1 spacetime has one time-like dimension ($x_0$, the direction of the octonion unit $e_0=1$) and seven space-like dimensions ($x_1,x_2,x_3$ ordinary space, $x_4,\dots,x_7$ auxiliary space). Two different objects describe this, and neither may be replaced by the other. (i) The tensor package xAct, used in section 12 of 7+1.nb exactly as in the author’s notebook Einstein-Lovelock-7+1-Nash.nb, declares the metric with \texttt{DefMetric[\{7, 1, 0\}, g71[-a, -b], CD, …]}: the list $\{7,1,0\}$ declares seven space-like dimensions, one time-like dimension and no degenerate dimension — the numbers of entries $+1$ and $-1$ of $\eta_{71}=\mathrm{diag}(-1,+1,\dots,+1)$ (\texttt{η71}), the same spacetime written in the sign convention of 4+4.nb (space-like $=+1$, time-like $=-1$). (ii) The matrix $\eta=-\eta_{71}=\mathrm{diag}(+1,-1,\dots,-1)$ (\texttt{η17}) is simply the squares of the octonion basis: $e_0^2=+1$ for the unit $e_0=1$ (time) and $e_i^2=-1$ for the seven imaginary units $e_1,\dots,e_7$ (space); more generally it is the real part of the product of two basis units, and it gives the quadratic form of the octonions:
\begin{gather*}
\eta_{AB}=\mathrm{Re}(e_Ae_B),\qquad e_0^2=+1,\qquad e_i^2=-1\ \ (i=1,\dots,7),\\ \mathrm{Re}(x\,x)=x_0^2-x_1^2-\cdots-x_7^2=\eta_{AB}\,x^Ax^B\qquad (x=x^Ae_A\in\mathbb O).
\end{gather*}
$\eta$ gives the signs of the Clifford relation $\Gamma_A\Gamma_B+\Gamma_B\Gamma_A=2\,\eta_{AB}\,1_{16}$ of the real 16×16 Dirac matrices: time squares to $+1_{16}$ and each of the seven space directions to $-1_{16}$, exactly like the octonion units (in the notation $Cl(p,q)$, with $p$ generators squaring to $+1$ and $q$ to $-1$, the algebra $Cl(1,7)\cong\mathbb R(16)$). It also gives the signs of the metric components: $g_{00}>0$ along the time-like direction, negative squares along the space-like directions. Every component computed by 7+1.nb follows from the component matrix of the metric; the declared list fixes for xAct the sign of the determinant of the metric, \texttt{SignDetOfMetric[g71]} $=-1$, which is the same in both sign conventions, since $\det \eta=\det \eta_{71}=-1$. 7+1.nb checks \texttt{SignatureOfMetric[g71]} $=\{7,1,0\}$, \texttt{SignDetOfMetric[g71]} $=-1$, $\mathrm{Re}(e_Ae_B)=\eta_{AB}$ and $\mathrm{Re}(x\,x)=\eta_{AB}x^Ax^B$.

\Needspace{12\baselineskip}\section*{2. Inputs from the earlier sections of 7+1.nb}\addcontentsline{toc}{section}{2. Inputs from the earlier sections of 7+1.nb}
{\footnotesize
\begin{longtable}{>{\raggedright\arraybackslash}p{\dimexpr 0.2500\linewidth-2\tabcolsep\relax}>{\raggedright\arraybackslash}p{\dimexpr 0.2500\linewidth-2\tabcolsep\relax}>{\raggedright\arraybackslash}p{\dimexpr 0.1000\linewidth-2\tabcolsep\relax}>{\raggedright\arraybackslash}p{\dimexpr 0.4000\linewidth-2\tabcolsep\relax}}
\toprule
\textbf{object} & \textbf{Mathematica} & \textbf{section} & \textbf{equation or value} \\
\midrule\endfirsthead
\toprule
\textbf{object} & \textbf{Mathematica} & \textbf{section} & \textbf{equation or value} \\
\midrule\endhead
\endfoot
\bottomrule\endlastfoot
stress tensor from the comoving energy & \texttt{feStress[W]} & 20.2 & $\varepsilon=W/(a^3b^4)$, $p_a=-\partial_aW/(3a^2b^4)$, $p_b=-\partial_bW/(4a^3b^3)$ \\
radion equation, radion mass & — & 20.3 & $\beta''+3H_a\beta'+4\beta'^2=\kappa(\frac\varepsilon6-\frac{p_a}2+\frac{p_b}3)+\frac\Lambda3$, $m_r^2=\kappa K_U/12$ \\
four-dimensional constants & \texttt{rIIMPl}, \texttt{rIIM8}, \texttt{rIIVp} & 20.3, 22.2 & $8\pi G_4=\kappa/V_\perp$, $V_\perp=M_{\rm Pl}^2/M_8^6$ \\
gas-dominated Einstein regime without stabilizer & \texttt{rIlate} & 21.3 & $\tau H_b\to\frac25$, $p_b/\varepsilon\to\frac14$ \\
gas share after capture & \texttt{rIIfrVals} & 22.1 & $\kappa\varepsilon/(3H_E^2)$ at $\tau=10^5$, $3\times10^5$ \\
radion decay into brane fields & \texttt{rIIgam} & 22.2 & $\sum\Gamma(\phi\to f\bar f)+\Gamma(\phi\to gg)$, $m_r=10$ GeV \\
radion decay into the bulk & \texttt{rIIgBulk} & 22.2 & $\sum_n\Gamma_n$ over the lattice, $2\mu_n<m_r$ \\
Jacobi’s $r_4(k)$ & \texttt{rIIr4} & 22.2 & $8\sigma(k)-32\sigma(k/4)$ for $k\le1.54\times10^6$ \\
late universe & \texttt{rIIH0g}, \texttt{rIIOL}, \texttt{rIIage}, \texttt{rIIgyr}, \texttt{rIImpc}, \texttt{rIIhbar}, \texttt{rIITcmb} & 22.4 & $H_0$ in GeV, $\Omega_\Lambda$, age 13.79 Gyr, conversion constants \\
Standard Model constants & \texttt{mHb}, \texttt{vEW}, \texttt{gFermi} & 17.3 & $m_h=125.20$ GeV, $v=(\sqrt2G_F)^{-1/2}$, $G_F=1.1663788\times10^{-5}$ GeV$^{-2}$ \\
uniform Kohn–Sham matter & \texttt{rho7}, \texttt{eps7}, \texttt{fermi7}, \texttt{gap7}, \texttt{gS0}, \texttt{gV0}, \texttt{kF0} & 8.7 & Hartree uniform matter in seven space dimensions, units $m=1$ \\*
\end{longtable}}
The formulas of these objects are written where they are used, in section 5. Before section 24 was appended, the evaluated notebook (sections 0–23, the summary of all checks numbered 24) was copied to \texttt{backups/\allowbreak{}7+1\_\allowbreak{}2026-10-27\_\allowbreak{}before\_\allowbreak{}dark.nb} (excluded from the repository by .gitignore); tools/verify\_7p1\_nb.py checks that every input and text cell and every output of sections 0–23 is unchanged, except one documented correction made with this stage: in the table of section 23.3 the cold-dark-matter density read "0.265"; it now reads "Omega\_c = 0.264", since $\Omega_ch^2/h^2=0.120/0.674^2=0.26416$.

\Needspace{12\baselineskip}\section*{3. Data and their sources}\addcontentsline{toc}{section}{3. Data and their sources}
{\footnotesize
\begin{longtable}{>{\raggedright\arraybackslash}p{\dimexpr 0.2200\linewidth-2\tabcolsep\relax}>{\raggedright\arraybackslash}p{\dimexpr 0.4000\linewidth-2\tabcolsep\relax}>{\raggedright\arraybackslash}p{\dimexpr 0.3800\linewidth-2\tabcolsep\relax}}
\toprule
\textbf{quantity} & \textbf{value used} & \textbf{source} \\
\midrule\endfirsthead
\toprule
\textbf{quantity} & \textbf{value used} & \textbf{source} \\
\midrule\endhead
\endfoot
\bottomrule\endlastfoot
$\Omega_ch^2$, $\Omega_bh^2$, $h$, $\Omega_m$ & 0.120, 0.0224, 0.674, 0.315 & Planck 2018 (Astron. Astrophys. 641 (2020) A6) \\
$T_{\rm CMB}$ & 2.7255 K & D. J. Fixsen, Astrophys. J. 707 (2009) 916 \\
$g_{*s0}$ & $2+\frac{21}4\cdot\frac4{11}=3.909$ & photons and three neutrino species after $e^\pm$ annihilation \\
DESI significances & $3.1\sigma$ (DESI + CMB); $2.8\sigma$–$4.2\sigma$ with supernovae & DESI DR2 Results II, arXiv:2503.14738 (abstract) \\
$w_0$, $w_a$ & $-0.753\pm0.056$, $-0.84^{+0.23}_{-0.20}$ (DESI + Planck PR4 + ACT DR6 + DES-Y5; $4.2\sigma$) & Garcia-Quintero et al., arXiv:2504.18464 (Tables II, III) \\
LZ limit & $2.2\times10^{-48}$ cm$^2$ at 40 GeV (90\% CL) & LZ Collaboration, arXiv:2410.17036 \\
self-interaction & $\sigma/m<1.25$ cm$^2$/g (68\% CL), $<0.7$ cm$^2$/g & Randall et al., ApJ 679 (2008) 1173 \\
warm dark matter & $m_{\rm WDM}>5.7$ keV (95\% CL) & Iršič et al., Phys. Rev. D 109 (2024) 043511 \\
decay-produced dark matter & about 7 keV for a thermal-relic bound of 3 keV & Ballesteros, Garcia, Pierre, JCAP 03 (2021) 101 (arXiv:2011.13458) \\
Pauli bound & 0.13 keV (95\% CL), $g=2$; scaling $g^{-1/4}$ (from their Eq. (9)) & Alvey et al., MNRAS 501 (2021) 1188 (arXiv:2010.03572) \\
X-ray lines & sterile neutrinos of 6–40 keV (qualitative use only) & Roach et al., Phys. Rev. D 107 (2023) 023009 (arXiv:2207.04572) \\
$M_D$ bounds & 5.9 to 11.2 TeV for $n=6$ to 2 & Particle Data Group, Phys. Rev. D 110 (2024) 030001 (ATLAS, 139 fb$^{-1}$) \\
$f_N$, $m_N$ & 0.30, 0.9383 GeV & Higgs–nucleon coupling factor (lattice-QCD sigma terms, common value); proton mass \\
$\alpha$ & 1/137.035999 & CODATA 2018 \\
$M_{\rm Pl}^{\rm (unreduced)}$ & $1.22091\times10^{19}$ GeV & $\sqrt{8\pi}\times2.435\times10^{18}$ GeV \\
$g_*$ in the freeze-in yield & 106.75 & all Standard Model particles relativistic; below the electroweak bosons $g_*=86.25$, so for $T_{\rm RH}<m_W$ the yield is understated by at most $(106.75/86.25)^{3/2}=1.4$ \\
1 GeV & $1.78266192\times10^{-24}$ g & CODATA 2018 \\
$\hbar c$ & $1.973269804\times10^{-14}$ GeV cm & CODATA 2018 \\
$M_8$, $m_r$ & 10 TeV, 10 GeV & choices of sections 22–23 (natural values) \\
$g_S$, $g_V$ (units $m=1$) & $-10^5$, $10^4$ & choices of section 8.7 (saturating uniform matter) \\*
\end{longtable}}
Software: Wolfram Language 15.0.1 (front end WolframNB.exe; the kernel configuration "Eval71" is a fresh kernel); Python 3.13.17 with numpy 2.5.3, scipy 1.18.1, sympy 1.14.0, mpmath 1.3.0.

\Needspace{12\baselineskip}\section*{4. Decisions}\addcontentsline{toc}{section}{4. Decisions}
{\footnotesize
\begin{longtable}{>{\raggedright\arraybackslash}p{\dimexpr 0.0500\linewidth-2\tabcolsep\relax}>{\raggedright\arraybackslash}p{\dimexpr 0.4700\linewidth-2\tabcolsep\relax}>{\raggedright\arraybackslash}p{\dimexpr 0.4800\linewidth-2\tabcolsep\relax}}
\toprule
\textbf{\#} & \textbf{decision} & \textbf{reason / equation} \\
\midrule\endfirsthead
\toprule
\textbf{\#} & \textbf{decision} & \textbf{reason / equation} \\
\midrule\endhead
\endfoot
\bottomrule\endlastfoot
1 & every statement labelled DERIVED, ESTIMATE or OPEN & the question mixes exact consequences of the field equations with model-dependent cosmology \\
2 & the dark-matter candidate is the lightest Kaluza–Klein mode $\chi=\psi\sqrt{V_\perp}$ (and, where relevant, the tower) & it is the lightest state; its couplings are fixed by the volume \\
3 & couplings of natural size at $M_8$: $\lambda_P=c_P/M_8^5$, $y=c_y/M_8^2$, $G_P=c_V/M_8^6$, $g_r=c_r/M_8^6$ & dimensional analysis of the brane and bulk operators \\
4 & decay rates of sterile neutrinos with the Majorana normalization & the standard formulas; a Dirac fermion decays half as fast (order one) \\
5 & stability criterion: lifetime longer than the age $t_0$ & the minimal requirement; X-ray limits are stronger and only strengthen the conclusions \\
6 & the overabundance of the primordial gas assuming the radion energy all became brane radiation & the most favourable case for the brane \\
7 & reheating after the dilution required by the moduli problem, instantaneous, at $T_{\rm RH}\le T_*$ & the modern regime starts below the normalcy temperature \\
8 & portal production estimated from Higgs decays with the Boltzmann factor $S(x)$ & a lower bound (other processes add); for the zero mode alone $S=1$ as an upper bound \\
9 & Kaluza–Klein sums in the continuum limit, checked against the lattice sums with $r_4(k)$ for $R=100$ & agreement to $10^{-9}$ \\
10 & vector portal: scaling estimate only & per-pair coupling of gravitational strength; the result is negligible by 17 orders of magnitude \\
11 & Pauli bound rescaled to $g=16$ & four brane Dirac spinors, two spins, particle and antiparticle \\
12 & dark energy from matter: the zero-temperature equation of state of the Kohn–Sham matter of section 8.7 & the self-coupled matter of the model with its saturation \\
13 & Casimir energy with periodic boundary conditions, dirac16complex and graviton only & the torus of the model; an estimate of the size of one contribution to $\Lambda_4$ \\
14 & the independent Python program uses its own quadratures, root finding and lattice sums; the coupled results come from tools/ref\_coupled.py & no code is shared with the notebook \\
15 & the rounding "0.265" in the table of section 23.3 corrected to "Omega\_c = 0.264" & $0.120/0.674^2=0.26416$; the only change to sections 0–23 \\*
\end{longtable}}
\Needspace{12\baselineskip}\section*{5. Derivations}\addcontentsline{toc}{section}{5. Derivations}
\textbf{The required abundance.} $\rho_{\rm DM}=\Omega_c\rho_c$ and $\rho_{\rm DM}=mn=mYs$ today give $mY=\Omega_ch^2(\rho_c/h^2)/s_0$ with $\rho_c=3H_0^2M_{\rm Pl}^2$ and $s_0=\frac{2\pi^2}{45}g_{*s0}T_{\rm CMB}^3$, $g_{*s0}=2+\frac78\cdot2\cdot3\cdot\frac4{11}=2+\frac{21}4\cdot\frac4{11}$ (neutrino temperature $(4/11)^{1/3}T$).

\textbf{The collider scale.} $M_{\rm Pl}^2=R^4M_D^6$ and $V_\perp=(2\pi R)^4=M_{\rm Pl}^2/M_8^6$ give $M_D^6=(2\pi)^4M_8^6$, $M_D=(2\pi)^{2/3}M_8$.

\textbf{A massive particle.} With $W=\omega=\sqrt{m^2+k_a^2/a^2+k_b^2/b^2}$: $\partial_aW=-k_a^2/(a^3\omega)$, $\partial_bW=-k_b^2/(b^3\omega)$, so
\begin{equation*}
\varepsilon a^3b^4=\omega,\qquad p_aa^3b^4=\frac{k_a^2}{3a^2\omega},\qquad p_ba^3b^4=\frac{k_b^2}{4b^2\omega},\qquad \frac{p_a}\varepsilon\Big|_{k_b=0}=\frac{k_a^2}{3(m^2a^2+k_a^2)} .
\end{equation*}
\textbf{The radion shift.} Static ($\beta'=\beta''=0$) and linear in $\beta$, the radion equation with the potential reads $m_r^2\beta=\kappa(\varepsilon/6-p_a/2+p_b/3)$ for the matter part; with $p_a=p_b=0$, $\kappa\varepsilon=\kappa\rho_{\rm DM}/V_\perp=\rho_{\rm DM}/M_{\rm Pl}^2=3H_0^2\Omega_c$:
\begin{equation*}
\delta\beta=\frac{\kappa\varepsilon}{6m_r^2}=\frac{H_0^2\Omega_c}{2m_r^2}.
\end{equation*}
\textbf{Regime III.} With $a^3=(\Omega_m/\Omega_\Lambda)\sinh^2(\frac32H_\infty t)$: $\Omega_ma^{-3}/\Omega_\Lambda=\sinh^{-2}(\frac32H_\infty t)$, and $\sinh^{-2}z=4e^{-2z}/(1-e^{-2z})^2\to4e^{-2z}$.

\textbf{Couplings of a level.} The brane value of the bulk field is $\psi(x,0)=\sum_n\chi_n(x)/\sqrt{V_\perp}$ (modes normalized on the torus), so an operator with $j$ factors of $\psi$ at the brane carries $V_\perp^{-j/2}$, and a bulk operator quartic in $\psi$, integrated over the torus, $V_\perp\cdot V_\perp^{-2}$:
\begin{equation*}
\lambda_{\rm eff}=\frac{\lambda_P}{V_\perp},\quad y_{\rm eff}=\frac{y}{\sqrt{V_\perp}},\quad G_{\rm eff}=\frac{G_P}{V_\perp},\quad G_r=\frac{g_r}{V_\perp},\qquad V_\perp=\frac{M_{\rm Pl}^2}{M_8^6}.
\end{equation*}
\textbf{Direct detection.} $\Phi^\dagger\Phi=(v+h)^2/2$ gives $-\lambda_{\rm eff}vh\bar\chi\chi$; with $-(f_Nm_N/v)h\bar NN$ and the $h$ propagator at zero momentum transfer, $\mathcal L=G_{\chi N}\bar\chi\chi\bar NN$, $G_{\chi N}=\lambda_{\rm eff}f_Nm_N/m_h^2$. In the non-relativistic limit $|\mathcal M|^2=G_{\chi N}^2(2m)^2(2m_N)^2$ and $\sigma=|\mathcal M|^2/(16\pi(m+m_N)^2)=G_{\chi N}^2\mu^2/\pi$.

\textbf{Mixing and decays.} In the basis $(\nu_L,\chi_L)\times(\nu_R,\chi_R)$ the zero mode and a neutrino have the mass matrix $\big(\begin{smallmatrix}m_\nu&\mu_P\\0&m\end{smallmatrix}\big)$ (section 18.2 of 7+1.nb); for $m\gg\mu_P$ the left-handed mixing angle is $\theta\approx\mu_P/m$. The heavy state then decays through the weak interactions of its $\nu_L$ admixture: $\Gamma(\chi\to3\nu)$ is the muon-decay rate $G_F^2m^5/(192\pi^3)$ times $\theta^2$ times two for the Majorana normalization (both charge-conjugate final states), and $\chi\to\nu\gamma$ is the one-loop magnetic transition (Pal and Wolfenstein), $\Gamma(\chi\to\nu\gamma)/\Gamma(\chi\to3\nu)=27\alpha/(8\pi)$. Since $\Gamma\propto\theta^2m^5=\mu_P^2m^3$, $\tau=K/(c_y^2m^3)$ and $m_{\max}=(K/(c_y^2t_0))^{1/3}$.

\textbf{The overabundance.} $\rho_{\rm rad}/s=\frac{\pi^2}{30}g_*T^4/(\frac{2\pi^2}{45}g_{*s}T^3)=\frac34T$ for $g_*=g_{*s}$; the gas, a fraction $f$ of the energy before the decay, is $f/(1-f)$ of the radiation after it; the dark matter needs $mYs/\rho_{\rm rad}=\frac43mY/T_{\rm RH}$.

\textbf{Freeze-in.} For $B\to\chi+\dots$ with $B$ in equilibrium ($n_\chi\ll n_\chi^{\rm eq}$), the collision term is $g_B\int\frac{d^3p}{(2\pi)^3}\frac{m_B}{E}\Gamma e^{-E/T}=\frac{g_B\Gamma m_B^2TK_1(m_B/T)}{2\pi^2}$. With $dt=-dT/(HT)$ (constant $g_{*s}$), $H=1.66\sqrt{g_*}T^2/M_{\rm Pl}^{\rm (unreduced)}$, $s=\frac{2\pi^2}{45}g_{*s}T^3$ and $x=m_B/T$:
\begin{align*}
Y&=\int_{T\le T_{\rm RH}}\frac{g_B\Gamma m_B^2TK_1(m_B/T)}{2\pi^2}\,\frac{dT}{H\,T\,s}=\frac{45\,g_B\Gamma M_{\rm Pl}^{\rm (unreduced)}}{4\pi^4\cdot1.66\sqrt{g_*}\,g_{*s}\,m_B^2}\int_{x_{\rm RH}}^\infty x^3K_1(x)\,dx, \\
\int_0^\infty x^3K_1(x)\,dx&=2^2\,\Gamma\big(\tfrac52\big)\Gamma\big(\tfrac32\big)=\frac{3\pi}2\quad\Rightarrow\quad Y=\frac{135\,g_B}{8\pi^3\cdot1.66\,g_*^{3/2}}\,\frac{M_{\rm Pl}^{\rm (unreduced)}\Gamma}{m_B^2}\,S(x_{\rm RH}),
\end{align*}
using $\int_0^\infty x^\mu K_\nu(x)dx=2^{\mu-1}\Gamma(\frac{1+\mu+\nu}2)\Gamma(\frac{1+\mu-\nu}2)$. \textbf{Higgs decays.} For $\mathcal L=-g\,h\bar\nu_L\chi_R+\text{h.c.}$, $g=y_{\rm eff}/\sqrt2$: $\sum|\mathcal M|^2=2g^2p_\nu\cdot p_\chi=g^2(m_h^2-\mu^2)$, $|\vec p|=(m_h^2-\mu^2)/(2m_h)$, so $\Gamma(h\to\bar\nu\chi)=y_{\rm eff}^2m_h(1-\mu^2/m_h^2)^2/(32\pi)$, and with the conjugate channel $y_{\rm eff}^2m_h/(16\pi)$ for $\mu\to0$. For $\mathcal L=-g\,h\bar\chi_1\chi_2$, $\sum|\mathcal M|^2=4g^2(p_1\cdot p_2-\mu_1\mu_2)=2g^2(m_h^2-(\mu_1+\mu_2)^2)$ and $|\vec p|=m_h\lambda^{1/2}(1,u_1^2,u_2^2)/2$:
\begin{equation*}
\Gamma(h\to\chi_1\bar\chi_2)=\frac{g^2m_h}{8\pi}\big(1-(u_1+u_2)^2\big)\lambda^{1/2}(1,u_1^2,u_2^2),\qquad \lambda(1,x_1,x_2)=\big(1-(\sqrt{x_1}+\sqrt{x_2})^2\big)\big(1-(\sqrt{x_1}-\sqrt{x_2})^2\big).
\end{equation*}
\textbf{Sums over the levels.} $\sum_{\vec n}F(|\vec n|)\approx\int d^4n\,F=2\pi^2\int_0^\infty n^3F(n)\,dn$ (the volume of the unit three-sphere is $2\pi^2$); with $|\vec n|=uR$, $R=m_h\ell_\perp/(2\pi)$:
\begin{align*}
\sum_n\Big(1-\frac{\mu_n^2}{m_h^2}\Big)^2&=2\pi^2R^4\int_0^1u^3(1-u^2)^2du=\frac{\pi^2R^4}{12}=\frac{V_\perp m_h^4}{192\pi^2}, \\
\sum_n\mu_n\Big(1-\frac{\mu_n^2}{m_h^2}\Big)^2&=2\pi^2R^4m_h\int_0^1u^4(1-u^2)^2du=\frac{16\pi^2R^4m_h}{315}=\frac{V_\perp m_h^5}{315\pi^2},
\end{align*}
and for the pairs $\sum_{n,n'}\to(V_\perp m_h^4/(8\pi^2))^2\iint u_1^3u_2^3(\cdots)du_1du_2$ from $dN=V_\perp\mu^3d\mu/(8\pi^2)$. \textbf{Vector portal.} Rate per pair $\sim G_{\rm eff}^2T^5$, pairs lighter than $T$ $\sim N(T)^2=(V_\perp T^4/(32\pi^2))^2$, $H\sim T^2/M_{\rm Pl}$: $\Gamma/H\sim c_V^2T^{11}V_\perp^2M_{\rm Pl}/M_{\rm Pl}^4=c_V^2T^{11}M_{\rm Pl}/M_8^{12}$ up to numbers, i.e. $(T/T_V)^{11}$ with $T_V^{11}=M_8^{12}/M_{\rm Pl}$.

\textbf{The equation of state of matter.} At zero temperature $E(N,V)=V\varepsilon(N/V)$ and $p=-\partial E/\partial V|_N=\rho\varepsilon'(\rho)-\varepsilon=\rho^2(\varepsilon/\rho)'$; the Hugenholtz–Van Hove theorem gives $\varepsilon'(\rho)=\varepsilon_F$, so $p=\varepsilon_F\rho-\varepsilon$ (checked in section 8.7 of 7+1.nb). A mixture of droplets at $\rho_0$ (volume fraction $\rho/\rho_0$) and empty space has the energy density $\rho\,\varepsilon_0/\rho_0$, lower than the uniform $\varepsilon(\rho)=\rho\,(\varepsilon/\rho)$ wherever $\varepsilon/\rho>\varepsilon_0/\rho_0$; the uniform phase with $p<0$ is therefore metastable ($dp/d\rho>0$) or unstable ($dp/d\rho<0$). \textbf{A condensate.} A Lorentz-invariant state has $\langle T_{\mu\nu}\rangle$ proportional to the only invariant tensor $g_{\mu\nu}$, hence $p=-\varepsilon$.

\textbf{The Casimir energy.} For a real scalar of mass $m$ on $\mathbb R^{3,1}\times T^4$ (periods $\ell_\perp$), the vacuum energy density per unit eight-volume is, after Poisson resummation over the Kaluza–Klein momenta and removal of the term $\vec n=0$ (the divergent flat-space part, absorbed into $\Lambda$),
\begin{align*}
\rho_8&=-\frac2{(2\pi)^4}\sum_{\vec n\ne0}\Big(\frac m{\ell_\perp|\vec n|}\Big)^4K_4(m\ell_\perp|\vec n|),\qquad \rho_{\rm Cas}^{(1)}=V_\perp\rho_8=-\frac2{\ell_\perp^4}\sum_{\vec n\ne0}\Big(\frac{m\ell_\perp}{2\pi|\vec n|}\Big)^4K_4(m\ell_\perp|\vec n|), \\
K_4(z)&\to\frac{\Gamma(4)}2\Big(\frac2z\Big)^4\ (z\to0):\qquad \rho_{\rm Cas}^{(1)}\to-\frac{\Gamma(4)}{\pi^4\ell_\perp^4}\sum_{\vec n\ne0}\frac1{|\vec n|^8}, \\
\sum_{k\ge1}\frac{r_4(k)}{k^s}&=8\sum_k\frac{\sigma(k)}{k^s}-32\sum_k\frac{\sigma(k)}{(4k)^s}=8(1-4^{1-s})\zeta(s)\zeta(s-1)\quad(s=4).
\end{align*}
A fermion loop contributes with the opposite sign and one real degree of freedom per on-shell state; dirac16complex has 16 such states, the eight-dimensional graviton $D(D-3)/2=20$.

\Needspace{12\baselineskip}\section*{6. Verification}\addcontentsline{toc}{section}{6. Verification}
\textbf{Evaluation of 7+1.nb}: Wolfram 15.0.1 for Microsoft Windows (64-bit) (July 2, 2026), a fresh kernel (front-end kernel configuration Eval71); finished Wed 28 Oct 2026 08:57:47; 283.883 s; message cells 0; output cells 554. The checks of section 24:

{\footnotesize
\begin{longtable}{>{\raggedright\arraybackslash}p{\dimexpr 0.0550\linewidth-2\tabcolsep\relax}>{\raggedright\arraybackslash}p{\dimexpr 0.7450\linewidth-2\tabcolsep\relax}>{\raggedright\arraybackslash}p{\dimexpr 0.0800\linewidth-2\tabcolsep\relax}>{\raggedright\arraybackslash}p{\dimexpr 0.1200\linewidth-2\tabcolsep\relax}}
\toprule
\textbf{\#} & \textbf{check} & \textbf{result} & \textbf{section} \\
\midrule\endfirsthead
\toprule
\textbf{\#} & \textbf{check} & \textbf{result} & \textbf{section} \\
\midrule\endhead
\endfoot
\bottomrule\endlastfoot
1 & \texttt{required abundance: m Y = Omega\_\allowbreak{}c h\textasciicircum{}2 (rho\_\allowbreak{}c/\allowbreak{}h\textasciicircum{}2)/\allowbreak{}s0 = 4.37 10\textasciicircum{}-10 GeV with rho\_\allowbreak{}c/\allowbreak{}h\textasciicircum{}2 = 1.053 10\textasciicircum{}-5 GeV/\allowbreak{}cm\textasciicircum{}3 and s0 = 2891 cm\textasciicircum{}-3 (to 10\textasciicircum{}-3)} & True & 24.1 \\
2 & \texttt{collider bound: with M\_\allowbreak{}Pl\textasciicircum{}2 = R\textasciicircum{}n M\_\allowbreak{}D\textasciicircum{}(n+2) and V\_\allowbreak{}perp = (2 Pi R)\textasciicircum{}4 = M\_\allowbreak{}Pl\textasciicircum{}2/\allowbreak{}M8\textasciicircum{}6, M\_\allowbreak{}D = (2 Pi)\textasciicircum{}(2/\allowbreak{}3) M8 = 34 TeV, above the strongest bound 11.2 TeV} & True & 24.1 \\
3 & \texttt{a massive dirac16complex particle: eps V == omega, pa V == k\_\allowbreak{}a\textasciicircum{}2/\allowbreak{}(3 a\textasciicircum{}2 omega), pb V == k\_\allowbreak{}b\textasciicircum{}2/\allowbreak{}(4 b\textasciicircum{}2 omega), omega = Sqrt[m\textasciicircum{}2 + k\_\allowbreak{}a\textasciicircum{}2/\allowbreak{}a\textasciicircum{}2 + k\_\allowbreak{}b\textasciicircum{}2/\allowbreak{}b\textasciicircum{}2]} & True & 24.2 \\
4 & \texttt{the lightest mode (k\_\allowbreak{}b = 0) is cold matter: pa/\allowbreak{}eps == k\_\allowbreak{}a\textasciicircum{}2/\allowbreak{}(3 (m\textasciicircum{}2 a\textasciicircum{}2 + k\_\allowbreak{}a\textasciicircum{}2)) -> 0, pb == 0, and eps \textasciitilde{} a\textasciicircum{}-3 at fixed b for k\_\allowbreak{}a = 0} & True & 24.2 \\
5 & \texttt{dark matter displaces the radion today by H0\textasciicircum{}2 Omega\_\allowbreak{}c/\allowbreak{}(2 m\_\allowbreak{}r\textasciicircum{}2) < 10\textasciicircum{}-80 (m\_\allowbreak{}r = 10 GeV): the four-dimensional Newton constant is constant} & True & 24.2 \\
6 & \texttt{regime III: Omega\_\allowbreak{}m a\textasciicircum{}-3/\allowbreak{}Omega\_\allowbreak{}Lambda == Csch[3 H t/\allowbreak{}2]\textasciicircum{}2 for a = (Omega\_\allowbreak{}m/\allowbreak{}Omega\_\allowbreak{}Lambda)\textasciicircum{}(1/\allowbreak{}3) Sinh[3 H t/\allowbreak{}2]\textasciicircum{}(2/\allowbreak{}3), and -> 4 Exp[-3 H t]} & True & 24.2 \\
7 & \texttt{couplings of the lightest mode: with V\_\allowbreak{}perp = M\_\allowbreak{}Pl\textasciicircum{}2/\allowbreak{}M8\textasciicircum{}6, lambda\_\allowbreak{}P/\allowbreak{}V\_\allowbreak{}perp == c\_\allowbreak{}P M8/\allowbreak{}M\_\allowbreak{}Pl\textasciicircum{}2, y/\allowbreak{}Sqrt[V\_\allowbreak{}perp] == c\_\allowbreak{}y M8/\allowbreak{}M\_\allowbreak{}Pl, G\_\allowbreak{}P/\allowbreak{}V\_\allowbreak{}perp == c\_\allowbreak{}V/\allowbreak{}M\_\allowbreak{}Pl\textasciicircum{}2, g\_\allowbreak{}r/\allowbreak{}V\_\allowbreak{}perp == c\_\allowbreak{}r/\allowbreak{}M\_\allowbreak{}Pl\textasciicircum{}2} & True & 24.3 \\
8 & \texttt{natural couplings make the lightest mode dark and collisionless: sigma\_\allowbreak{}SI(c\_\allowbreak{}P = 1, m = 40 GeV) < 10\textasciicircum{}-50 x (LZ limit 2.2 10\textasciicircum{}-48 cm\textasciicircum{}2), sigma/\allowbreak{}m(c\_\allowbreak{}r = 1, m = 1 GeV) < 10\textasciicircum{}-70 cm\textasciicircum{}2/\allowbreak{}g (Bullet Cluster: 0.7 cm\textasciicircum{}2/\allowbreak{}g)} & True & 24.3 \\
9 & \texttt{stability through the Yukawa portal: for c\_\allowbreak{}y = 1 the lifetime chi -> 3 nu equals the age of the universe at m = 40 MeV; it scales as 1/\allowbreak{}(c\_\allowbreak{}y\textasciicircum{}2 m\textasciicircum{}3), and Gamma(nu gamma)/\allowbreak{}Gamma(3 nu) == 27 alpha/\allowbreak{}(8 Pi)} & True & 24.3 \\
10 & \texttt{primordial gas: overabundance O = (f/\allowbreak{}(1 - f)) T\_\allowbreak{}RH/\allowbreak{}((4/\allowbreak{}3) m Y) = 1.5 10\textasciicircum{}7 T\_\allowbreak{}RH/\allowbreak{}GeV (6 10\textasciicircum{}4 at T\_\allowbreak{}RH = 4 MeV)} & True & 24.4 \\
11 & \texttt{the radion puts more than 1 - 10\textasciicircum{}-12 of its decays into the bulk (Kaluza-Klein pairs of dirac16complex)} & True & 24.4 \\
12 & \texttt{Boltzmann factor of production below T\_\allowbreak{}RH: Int\_\allowbreak{}0\textasciicircum{}Infinity t\textasciicircum{}3 K\_\allowbreak{}1(t) dt == 3 Pi/\allowbreak{}2 (exact), S(0) = 1, S(m\_\allowbreak{}h/\allowbreak{}T*) < 10\textasciicircum{}-2 with T* = (M8\textasciicircum{}6/\allowbreak{}M\_\allowbreak{}Pl)\textasciicircum{}(1/\allowbreak{}5) = 13.3 GeV} & True & 24.4 \\
13 & \texttt{zero mode through the Yukawa portal: the abundance (\textasciitilde{} c\_\allowbreak{}y\textasciicircum{}2 m) and the stability (lifetime \textasciitilde{} 1/\allowbreak{}(c\_\allowbreak{}y\textasciicircum{}2 m\textasciicircum{}3) > age) are compatible only for m < m*, and m* < 5.7 keV (Lyman-alpha) for N = 1 and N = 12 channels even with S = 1} & True & 24.4 \\
14 & \texttt{Kaluza-Klein tower: Sum\_\allowbreak{}n (1 - mu\_\allowbreak{}n\textasciicircum{}2/\allowbreak{}m\_\allowbreak{}h\textasciicircum{}2)\textasciicircum{}2 -> V\_\allowbreak{}perp m\_\allowbreak{}h\textasciicircum{}4/\allowbreak{}(192 Pi\textasciicircum{}2) and Sum\_\allowbreak{}n mu\_\allowbreak{}n (1 - mu\_\allowbreak{}n\textasciicircum{}2/\allowbreak{}m\_\allowbreak{}h\textasciicircum{}2)\textasciicircum{}2 -> V\_\allowbreak{}perp m\_\allowbreak{}h\textasciicircum{}5/\allowbreak{}(315 Pi\textasciicircum{}2) (d\textasciicircum{}4n = 2 Pi\textasciicircum{}2 |n|\textasciicircum{}3 d|n|: exact integrals); the lattice sums with r\_\allowbreak{}4(k) for R = 100 equal Pi\textasciicircum{}2 R\textasciicircum{}4/\allowbreak{}12 and 16 Pi\textasciicircum{}2 R\textasciicircum{}4/\allowbreak{}315 to 10\textasciicircum{}-8} & True & 24.4 \\
15 & \texttt{Kaluza-Klein tower through the Yukawa portal: 7.7 10\textasciicircum{}17 states below m\_\allowbreak{}h with mean mass (192/\allowbreak{}315) m\_\allowbreak{}h; at T\_\allowbreak{}RH = T* it carries more energy than all the dark matter unless c\_\allowbreak{}y < 2 10\textasciicircum{}-5 (N = 1)} & True & 24.4 \\
16 & \texttt{Higgs portal: the zero mode gets Y < 10\textasciicircum{}-40 (c\_\allowbreak{}P = 1); the pairs of Kaluza-Klein levels (lambda(1, u1\textasciicircum{}2, u2\textasciicircum{}2) == (1 - (u1 + u2)\textasciicircum{}2)(1 - (u1 - u2)\textasciicircum{}2)) give the observed abundance only for c\_\allowbreak{}P between 10 and 20 (T\_\allowbreak{}RH >> m\_\allowbreak{}h, S -> 1) or between 100 and 200 (T\_\allowbreak{}RH = T*)} & True & 24.4 \\
17 & \texttt{vector portal: gravitational strength per pair; the production rate of all pairs lighter than T divided by H \textasciitilde{} (T/\allowbreak{}T\_\allowbreak{}V)\textasciicircum{}11 with T\_\allowbreak{}V = (M8\textasciicircum{}12/\allowbreak{}M\_\allowbreak{}Pl)\textasciicircum{}(1/\allowbreak{}11) = 490 GeV is below 10\textasciicircum{}-17 at T*} & True & 24.4 \\
18 & \texttt{Kohn-Sham matter of section 8: p > 0 at low density, p < 0 below saturation, p = 0 at kF0, a spinodal region (dp/\allowbreak{}dkF < 0) in between, and w = p/\allowbreak{}eps >= -0.06 everywhere} & True & 24.5 \\
19 & \texttt{Epstein sum over Z\textasciicircum{}4 without 0: Sum |n|\textasciicircum{}-8 == 8 (1 - 4\textasciicircum{}-3) zeta(4) zeta(3), compared with the shell sum Sum r\_\allowbreak{}4(k)/\allowbreak{}k\textasciicircum{}4 to k = 20000 (tail < 3 10\textasciicircum{}-8)} & True & 24.5 \\
20 & \texttt{Casimir energy of dirac16complex on the auxiliary torus: massless 16 Gamma(4) Z/\allowbreak{}(Pi\textasciicircum{}4 l\textasciicircum{}4) > 10\textasciicircum{}34 rho\_\allowbreak{}Lambda, the massive sum -> the massless value for m l -> 0 (to 10\textasciicircum{}-4), and equals rho\_\allowbreak{}Lambda at m l = 92.2} & True & 24.5 \\*
\end{longtable}}
\textbf{The independent Python program} tools/ref\_dark.py — its boolean checks:

{\footnotesize
\begin{longtable}{>{\raggedright\arraybackslash}p{\dimexpr 0.0550\linewidth-2\tabcolsep\relax}>{\raggedright\arraybackslash}p{\dimexpr 0.8250\linewidth-2\tabcolsep\relax}>{\raggedright\arraybackslash}p{\dimexpr 0.1200\linewidth-2\tabcolsep\relax}}
\toprule
\textbf{\#} & \textbf{check} & \textbf{result} \\
\midrule\endfirsthead
\toprule
\textbf{\#} & \textbf{check} & \textbf{result} \\
\midrule\endhead
\endfoot
\bottomrule\endlastfoot
1 & \texttt{required abundance m Y = 4.37e-10 GeV (rho\_\allowbreak{}c/\allowbreak{}h\textasciicircum{}2 = 1.053e-5 GeV/\allowbreak{}cm\textasciicircum{}3, s0 = 2891/\allowbreak{}cm\textasciicircum{}3, to 1e-3)} & True \\
2 & \texttt{M\_\allowbreak{}Pl\textasciicircum{}2 = R\textasciicircum{}4 M\_\allowbreak{}D\textasciicircum{}6 with V = (2 pi R)\textasciicircum{}4 = M\_\allowbreak{}Pl\textasciicircum{}2/\allowbreak{}M8\textasciicircum{}6 gives M\_\allowbreak{}D = (2 pi)\textasciicircum{}(2/\allowbreak{}3) M8 = 34 TeV > 11.2 TeV (strongest collider bound)} & True \\
3 & \texttt{massive particle: eps V = omega, pa V = k\_\allowbreak{}a\textasciicircum{}2/\allowbreak{}(3 a\textasciicircum{}2 omega), pb V = k\_\allowbreak{}b\textasciicircum{}2/\allowbreak{}(4 b\textasciicircum{}2 omega)} & True \\
4 & \texttt{lightest mode (k\_\allowbreak{}b = 0): pa/\allowbreak{}eps = k\_\allowbreak{}a\textasciicircum{}2/\allowbreak{}(3 (m\textasciicircum{}2 a\textasciicircum{}2 + k\_\allowbreak{}a\textasciicircum{}2)), pb = 0, eps \textasciitilde{} a\textasciicircum{}-3 at fixed b for k\_\allowbreak{}a = 0} & True \\
5 & \texttt{dark matter leaves G4 unchanged: displacement < 1e-80} & True \\
6 & \texttt{regime III: Omega\_\allowbreak{}m a\textasciicircum{}-3/\allowbreak{}Omega\_\allowbreak{}L = csch\textasciicircum{}2(3 H t/\allowbreak{}2) for a = (Omega\_\allowbreak{}m/\allowbreak{}Omega\_\allowbreak{}L)\textasciicircum{}(1/\allowbreak{}3) sinh\textasciicircum{}(2/\allowbreak{}3)(3 H t/\allowbreak{}2), -> 4 exp(-3 H t)} & True \\
7 & \texttt{zero-mode couplings: lambda\_\allowbreak{}P/\allowbreak{}V = c\_\allowbreak{}P M8/\allowbreak{}M\_\allowbreak{}Pl\textasciicircum{}2, y/\allowbreak{}sqrt(V) = c\_\allowbreak{}y M8/\allowbreak{}M\_\allowbreak{}Pl, G\_\allowbreak{}P/\allowbreak{}V = c\_\allowbreak{}V/\allowbreak{}M\_\allowbreak{}Pl\textasciicircum{}2, g\_\allowbreak{}r/\allowbreak{}V = c\_\allowbreak{}r/\allowbreak{}M\_\allowbreak{}Pl\textasciicircum{}2} & True \\
8 & \texttt{dark and collisionless: sigma\_\allowbreak{}SI < 1e-50 x 2.2e-48 cm\textasciicircum{}2, sigma/\allowbreak{}m < 1e-70 cm\textasciicircum{}2/\allowbreak{}g} & True \\
9 & \texttt{stability: lifetime = age at m = 40 MeV (c\_\allowbreak{}y = 1); lifetime \textasciitilde{} 1/\allowbreak{}(c\_\allowbreak{}y\textasciicircum{}2 m\textasciicircum{}3); Gamma(nu gamma)/\allowbreak{}Gamma(3 nu) = 27 alpha/\allowbreak{}(8 pi)} & True \\
10 & \texttt{primordial gas overabundance 1.5e7 T\_\allowbreak{}RH/\allowbreak{}GeV (6e4 at 4 MeV)} & True \\
11 & \texttt{radion decays: more than 1 - 1e-12 into the bulk} & True \\
12 & \texttt{Boltzmann factor: int\_\allowbreak{}0\textasciicircum{}oo t\textasciicircum{}3 K\_\allowbreak{}1(t) dt = 3 pi/\allowbreak{}2 (sympy), S(0) = 1 (mpmath), S(m\_\allowbreak{}h/\allowbreak{}T*) < 1e-2, 13 < T* < 14 GeV} & True \\
13 & \texttt{zero mode, Yukawa portal: abundance and stability compatible only below m* < 5.7 keV (N = 1, 12)} & True \\
14 & \texttt{Kaluza-Klein tower: int 2 pi\textasciicircum{}2 u\textasciicircum{}3 (1-u\textasciicircum{}2)\textasciicircum{}2 = pi\textasciicircum{}2/\allowbreak{}12, int 2 pi\textasciicircum{}2 u\textasciicircum{}4 (1-u\textasciicircum{}2)\textasciicircum{}2 = 16 pi\textasciicircum{}2/\allowbreak{}315 (sympy); lattice sums agree to 1e-8 (R = 100)} & True \\
15 & \texttt{tower through the Yukawa portal: 7.7e17 states below m\_\allowbreak{}h, mean mass (192/\allowbreak{}315) m\_\allowbreak{}h, c\_\allowbreak{}y bound < 2e-5 at T*} & True \\
16 & \texttt{Higgs portal: zero mode Y < 1e-40; lambda(1, u1\textasciicircum{}2, u2\textasciicircum{}2) = (1 - (u1+u2)\textasciicircum{}2)(1 - (u1-u2)\textasciicircum{}2); pairs need 10 < c\_\allowbreak{}P < 20 (S = 1), 100 < c\_\allowbreak{}P < 200 (T*)} & True \\
17 & \texttt{vector portal: 480 < T\_\allowbreak{}V < 500 GeV, (T*/\allowbreak{}T\_\allowbreak{}V)\textasciicircum{}11 < 1e-17} & True \\
18 & \texttt{Kohn-Sham matter: p > 0 at low density, p < 0 just below kF0, p = 0 at kF0, a spinodal region, w >= -0.06} & True \\
19 & \texttt{Epstein sum: Sum over Z\textasciicircum{}4\textbackslash{}\allowbreak{}\{0\} of |n|\textasciicircum{}-8 = 8 (1 - 4\textasciicircum{}-3) zeta(4) zeta(3) = shell sum to k = 20000 (tail < 3e-8)} & True \\
20 & \texttt{Casimir: massless > 1e34 rho\_\allowbreak{}Lambda; massive sum -> massless (m l = 1e-3, to 1e-4); = rho\_\allowbreak{}Lambda at m l = 92.2} & True \\*
\end{longtable}}
\textbf{Numbers printed by both programs} (largest relative difference per group, and the tolerance used by tools/verify\_7p1\_nb.py):

{\footnotesize
\begin{longtable}{>{\raggedright\arraybackslash}p{\dimexpr 0.5200\linewidth-2\tabcolsep\relax}>{\raggedright\arraybackslash}p{\dimexpr 0.1000\linewidth-2\tabcolsep\relax}>{\raggedright\arraybackslash}p{\dimexpr 0.2200\linewidth-2\tabcolsep\relax}>{\raggedright\arraybackslash}p{\dimexpr 0.1600\linewidth-2\tabcolsep\relax}}
\toprule
\textbf{quantity} & \textbf{values} & \textbf{max. relative difference} & \textbf{tolerance} \\
\midrule\endfirsthead
\toprule
\textbf{quantity} & \textbf{values} & \textbf{max. relative difference} & \textbf{tolerance} \\
\midrule\endhead
\endfoot
\bottomrule\endlastfoot
$\Omega_c$, $\rho_c/h^2$, $s_0$, $mY$ & 4 & $2.4\times10^{-16}$ & $1\times10^{-12}$ \\
$M_D$ & 1 & 0 & $1\times10^{-14}$ \\
Pauli bound for $g=16$ & 1 & 0 & $1\times10^{-14}$ \\
radion displacement by dark matter & 1 & $4.9\times10^{-16}$ & $1\times10^{-12}$ \\
gas share after capture (from the coupled solution) & 2 & $1.5\times10^{-10}$ & $1\times10^{-8}$ \\
$\tau H_b$ without stabilizer (gas-dominated) & 1 & $1.4\times10^{-7}$ & $1\times10^{-5}$ \\
$\lambda_{\rm eff}$, $y_{\rm eff}$, $G_r$ & 3 & 0 & $1\times10^{-14}$ \\
$\sigma_{\rm SI}$ & 1 & $1.3\times10^{-16}$ & $1\times10^{-12}$ \\
$\sigma/m$; $c_r$ allowed & 2 & $1.2\times10^{-16}$ & $1\times10^{-12}$ \\
$\mu_P$; lifetimes & 5 & $2.8\times10^{-16}$ & $1\times10^{-12}$ \\
largest stable $m$ ($c_y=1$) & 1 & $3.3\times10^{-12}$ & $1\times10^{-9}$ \\
gas share $f$, overabundance $O$ & 3 & $2.2\times10^{-12}$ & $1\times10^{-8}$ \\
radion decays: bulk fraction, brane share & 2 & $3.3\times10^{-16}$ & $1\times10^{-9}$ \\
$T_*$, $S(m_h/T_*)$ & 2 & $1.1\times10^{-15}$ & $1\times10^{-9}$ \\
zero mode: $Y$ (Yukawa, Higgs portal) & 2 & $3.2\times10^{-16}$ & $1\times10^{-12}$ \\
$m_*$ ($N=1$, 12) & 2 & $8.2\times10^{-16}$ & $1\times10^{-9}$ \\
$c_y$ needed and allowed & 4 & $6.9\times10^{-16}$ & $1\times10^{-9}$ \\
Kaluza–Klein tower (Yukawa portal) & 5 & $5.1\times10^{-16}$ & $1\times10^{-9}$ \\
Kaluza–Klein lifetime; largest stable $m$ ($c_y=10^{-5}$) & 2 & $6.6\times10^{-16}$ & $1\times10^{-9}$ \\
Higgs-portal pairs & 6 & $7.5\times10^{-14}$ & $1\times10^{-8}$ \\
Higgs-portal pairs: fraction of the dark matter & 1 & $7.5\times10^{-14}$ & $1\times10^{-8}$ \\
vector portal: $T_V$, $(T_*/T_V)^{11}$ & 2 & 0 & $1\times10^{-12}$ \\
dependence on $T_{\rm RH}$ & 16 & $3.8\times10^{-14}$ & $1\times10^{-8}$ \\
Kohn–Sham matter: $w_{\min}$, its $k_F$, $k_{F0}$ & 3 & 0 & $1\times10^{-9}$ \\
Kohn–Sham matter: $(k_F,p,w)$ table & 21 & 0 & $1\times10^{-9}$ \\
$\rho_\Lambda$ & 1 & $5.8\times10^{-16}$ & $1\times10^{-12}$ \\
Casimir energies (massless) & 2 & $2.2\times10^{-16}$ & $1\times10^{-12}$ \\
$m\ell_\perp$ and $m$ where the Casimir energy equals $\rho_\Lambda$ & 2 & $7.7\times10^{-16}$ & $1\times10^{-9}$ \\
Casimir energy at $m\ell_\perp=1,10,30,60$ & 4 & $1.7\times10^{-16}$ & $1\times10^{-9}$ \\*
\end{longtable}}
\textbf{SHA-256 of the files} (at the time this document was generated):

\begin{Verbatim}[fontsize=\small,frame=single,framesep=2mm]
7+1.nb                             41d252f9852f210c228f45e1ac956d2d7ef26faeb0062a5880280e1aae0ef
    ...fc7
tools/make_7p1_nb.py               996d26c9552cd44902560af16a1683fa71457261f18fefaeddae87bd99a2b
    ...144
tools/nb_part_gravity.py           c4caf2ad5c932e8f603341edade5f4b98fa784feb05906c051eab7c884cd6
    ...8f8
tools/nb_part_ks_gravity.py        a656c0ad040b4576b75b17e029545109e531c8e656ad99bfdb5cba946e967
    ...772
tools/nb_part_sm.py                43e92beed57453a4b57931de6d603ac6c218efba6e4005133b69b8054a975
    ...c5b
tools/nb_part_coupled.py           e630521d64132e71afc861114f699eb06454557cf14a57ad8ec1eef96247a
    ...ae0
tools/nb_part_dark.py              10ad8fb8ae25fa07bbee758c653995acd4e2c80733059a8b3648bd1593e46
    ...db0
tools/nb_part_torus_brane.py       fb904ff8d85b904c55b3b2a8f60fb0cb047f1ccf1b98db3364092f2bae0b2
    ...162
tools/nb_part_radion.py            53c5604707a3e22228c840ee0a2d2125616f9410f75a5535041dfaf18f497
    ...c00
tools/lovelock71.py                d93ecbd0f82adb972218d59dfe32999a4af944945ca9758607f3a633e1af9
    ...5a1
tools/ref_gravity.py               f07cf8adaff585af89c87d19532f1aaa7db606ad88ae90483b2d1f41341ed
    ...f36
tools/ref_ks_gravity.py            b6488ef24b1c16aacd8ff6cf6a5a1ba3f316e082449ad97480022d1a26383
    ...27f
tools/ref_sm.py                    f3f80806893f42bbf166e1f1de896ffd611dbfe847685da8e3ee367e9cde6
    ...aba
tools/ref_coupled.py               695dbd9375753653472d9c37502f0067c4a64fe01a64ba8c05e8b2c5da85a
    ...a58
tools/ref_dark.py                  1600e6b57ab0383b4f2b664e2796d3864c9afa77a9ed59f007f66e1139453
    ...ff5
tools/ref_torus_brane.py           1f623656ab0f5c8ee43cb310f7bfd580cbba5a7385da2963adaae4d2f8bda
    ...af7
tools/ref_radion.py                e9d47176b9f2d2232a84b3b82188bd3b0a19cde3dbe1ba013c1fe5abff634
    ...1ce
tools/ref71.py                     1a7a06eb0add3f459611776dc9c474a5bb959d947a9db949db4612b1ab943
    ...4e0
tools/ref44.py                     ab100c3ebcc055fbe59c132a02f80ecb17630e934e99d59f88e919ec8064f
    ...af5
\end{Verbatim}
\textbf{Output of tools/ref\_dark.py (independent checks of section 24; exit code 0)} (complete):

\begin{Verbatim}[fontsize=\small,frame=single,framesep=2mm]
NUM Omega_c, rho_c/h^2 (GeV/cm^3), s0 (1/cm^3), m Y (GeV) : (0.264156592027754,
    ...1.0533924943202557e-05, 2891.2844014437223, 4.372005025009337e-10)
NUM M_D (GeV) = (2 pi)^(2/3) M8 : 34050.219214767545
NUM Pauli bound for g = 16 internal states, 0.13 keV (2/16)^(1/4) (keV) : 0.07729846247517687
required abundance m Y = 4.37e-10 GeV (rho_c/h^2 = 1.053e-5 GeV/cm^3, s0 = 2891/cm^3, to 1e-3)
    ...: True
M_Pl^2 = R^4 M_D^6 with V = (2 pi R)^4 = M_Pl^2/M8^6 gives M_D = (2 pi)^(2/3) M8 = 34 TeV >
    ...11.2 TeV (strongest collider bound) : True
massive particle: eps V = omega, pa V = k_a^2/(3 a^2 omega), pb V = k_b^2/(4 b^2 omega) : True
lightest mode (k_b = 0): pa/eps = k_a^2/(3 (m^2 a^2 + k_a^2)), pb = 0, eps ~ a^-3 at fixed b
    ...for k_a = 0 : True
NUM radion displacement by dark matter today H0^2 Omega_c/(2 m_r^2) (m_r = 10 GeV) :
    ...2.730119736367085e-87
dark matter leaves G4 unchanged: displacement < 1e-80 : True
regime III: Omega_m a^-3/Omega_L = csch^2(3 H t/2) for a = (Omega_m/Omega_L)^(1/3) sinh^(2/3)(3
    ...H t/2), -> 4 exp(-3 H t) : True
NUM from tools/ref_coupled.py: gas share at tau = 1e5, 3e5; Gamma_brane, Gamma_bulk (GeV); tau
    ...Hb at Ha = 1e-14 : (0.008595792799709226, 0.008728032079226138, 4.572168576618335e-37,
    ...5.9991693514127285e-24, 0.3999952443921748)
zero-mode couplings: lambda_P/V = c_P M8/M_Pl^2, y/sqrt(V) = c_y M8/M_Pl, G_P/V = c_V/M_Pl^2,
    ...g_r/V = c_r/M_Pl^2 : True
NUM lambda_eff, y_eff, G_r per unit coupling (GeV^-1, 1, GeV^-2) : (1.68656105983497e-33,
    ...4.106776180698152e-15, 1.68656105983497e-37)
NUM sigma_SI (cm^2) at c_P = 1, m = 40 GeV : 9.556076962877839e-104
NUM sigma/m (cm^2/g) at c_r = 1, m = 1 GeV; c_r allowed by the Bullet Cluster at m = 1 GeV :
    ...(4.9442250036435816e-79, 1.1898710736040751e+39)
NUM mu_P (GeV) at c_y = 1; lifetimes chi -> nu gamma, chi -> 3 nu (s) at c_y = 1, m = 1 GeV and
    ...10 keV : (7.150044241252333e-13, 3593359688030266.0, 28170199696474.52,
    ...3.5933596880302647e+30, 2.8170199696474517e+28)
NUM largest m (GeV) with lifetime chi -> 3 nu > age for c_y = 1 : 0.04015133785290378
dark and collisionless: sigma_SI < 1e-50 x 2.2e-48 cm^2, sigma/m < 1e-70 cm^2/g : True
stability: lifetime = age at m = 40 MeV (c_y = 1); lifetime ~ 1/(c_y^2 m^3); Gamma(nu
    ...gamma)/Gamma(3 nu) = 27 alpha/(8 pi) : True
NUM gas share after capture f, overabundance O at T_RH = 1 GeV and 4 MeV :
    ...(0.008728032079226138, 15104422.311749911, 60417.689246999646)
NUM fraction of radion decays into the bulk; brane share Gamma_brane/Gamma_bulk :
    ...(0.9999999999999238, 7.621336069703795e-14)
primordial gas overabundance 1.5e7 T_RH/GeV (6e4 at 4 MeV) : True
radion decays: more than 1 - 1e-12 into the bulk : True
NUM T* (GeV), S(m_h/T*) : (13.264784746562139, 0.0078461320088985)
Boltzmann factor: int_0^oo t^3 K_1(t) dt = 3 pi/2 (sympy), S(0) = 1 (mpmath), S(m_h/T*) < 1e-2,
    ...13 < T* < 14 GeV : True
NUM zero mode: Yukawa-portal Y at c_y = 1 (N = 1, S = 1), Higgs-portal Y at c_P = 1 (S = 1) :
    ...(9.726257996829745e-18, 1.9889471111127847e-49)
NUM m* (GeV) for N = 1 and N = 12 : (1.2000037341243723e-06, 4.156934873551575e-06)
NUM c_y needed and allowed at m = 10 keV, 10 MeV : (2120154.126594499, 254419.2868832596,
    ...67045.16030643586, 8.045444272269018)
zero mode, Yukawa portal: abundance and stability compatible only below m* < 5.7 keV (N = 1,
    ...12) : True
NUM lattice sums for R = 100: Sum r_4(k)(1 - k/R^2)^2 / (pi^2 R^4/12) - 1, Sum r_4(k) sqrt(k)/R
    ...(1 - k/R^2)^2 / (16 pi^2 R^4/315) - 1 : (5.981171113944583e-10, 9.233644959749654e-10)
Kaluza-Klein tower: int 2 pi^2 u^3 (1-u^2)^2 = pi^2/12, int 2 pi^2 u^4 (1-u^2)^2 = 16 pi^2/315
    ...(sympy); lattice sums agree to 1e-8 (R = 100) : True
NUM Kaluza-Klein tower below m_h: number, mass sum (GeV), Y_tower and mass per entropy (GeV) at
    ...c_y = 1 (N = 1, S = 1); c_y bound at T* : (7.688015893943899e+17, 5.866907976666065e+19,
    ...7.477562606822603, 570.6306062471253, 9.881775656983552e-06)
NUM tau(chi_n -> 3 nu) (s) at mu_n = 10 GeV, c_y = 1; largest m (GeV) with tau(chi -> 3 nu) >
    ...age for c_y = 1e-5 : (28170199696.47452, 86.50343512177517)
tower through the Yukawa portal: 7.7e17 states below m_h, mean mass (192/315) m_h, c_y bound <
    ...2e-5 at T* : True
NUM Higgs portal pairs: I_hp, I_hp^E, number of pairs, mass per entropy (GeV) at c_P = 1 (N =
    ...4, S = 1); c_P needed for S = 1 and at T* : (9.538467456842521e-05,
    ...7.380786601565824e-05, 3.40448189100676e+38, 2.5028869407639637e-12, 13.216598897717843,
    ...149.20792978255014)
NUM Higgs portal pairs at c_P = 1, T_RH = T*: fraction of the dark matter :
    ...4.491756351661614e-05
Higgs portal: zero mode Y < 1e-40; lambda(1, u1^2, u2^2) = (1 - (u1+u2)^2)(1 - (u1-u2)^2);
    ...pairs need 10 < c_P < 20 (S = 1), 100 < c_P < 200 (T*) : True
NUM vector portal: T_V (GeV), (T*/T_V)^11 : (492.19416145801307, 5.4475502039269596e-18)
vector portal: 480 < T_V < 500 GeV, (T*/T_V)^11 < 1e-17 : True
NUM T_RH (GeV) = 5, 8, 10, T*: S(m_h/T_RH), c_y bound of the tower (N = 1), c_P needed by the
    ...Higgs-portal pairs (N = 4) : [(5.0, 1.248895314095822e-08, 0.007832488705773653,
    ...118265.1241437344), (8.0, 4.9440196437621006e-05, 0.00012448662751176695,
    ...1879.6613707290032), (10.0, 0.0006775813090297954, 3.3626532542268916e-05,
    ...507.73722057246755), (13.264784746562139, 0.0078461320088985, 9.881775656983552e-06,
    ...149.20792978255014)]
NUM w_min of Kohn-Sham matter and its kF (units m = 1); saturation kF0 : (-0.05801598575399315,
    ...0.67, 0.7566361507014929)
Kohn-Sham matter: p > 0 at low density, p < 0 just below kF0, p = 0 at kF0, a spinodal region,
    ...w >= -0.06 : True
Epstein sum: Sum over Z^4\{0} of |n|^-8 = 8 (1 - 4^-3) zeta(4) zeta(3) = shell sum to k = 20000
    ...(tail < 3e-8) : True
NUM rho_Lambda (GeV^4) : 2.518265996152132e-47
NUM Casimir energy / rho_Lambda: dirac16complex (m = 0), graviton (20 dof) :
    ...(6.762461109529047e+34, -8.45307638691131e+34)
NUM m l_perp and m (GeV) at which the dirac16complex Casimir energy equals rho_Lambda :
    ...(92.22624853874979, 0.05910239233816993)
NUM kF, p, w at kF = 0.1, 0.3, 0.5, 0.6, 0.67, 0.7, 0.75 (units m = 1) : [(0.1,
    ...1.0815057495004936e-14, 0.0011018832731922797), (0.3, 1.8769661856714918e-10,
    ...0.00849359466975192), (0.5, -4.3081067785473656e-10, -0.0005296332340018398), (0.6,
    ...-1.0277578600831483e-07, -0.035970422554837775), (0.67, -3.4548242506780997e-07,
    ...-0.05801598575399315), (0.7, -4.109652096104839e-07, -0.05166870169668898), (0.75,
    ...-1.0713521089577265e-07, -0.008446078650216108)]
NUM dirac16complex Casimir energy / rho_Lambda at m l_perp = 1, 10, 30, 60 :
    ...[6.05689541888988e+34, 4.1996195140622416e+32, 2.4693584837075517e+25,
    ...23006760008666.074]
Casimir: massless > 1e34 rho_Lambda; massive sum -> massless (m l = 1e-3, to 1e-4); =
    ...rho_Lambda at m l = 92.2 : True
boolean checks: 20   failed: 0
VERDICT: ALL CHECKS PASSED
\end{Verbatim}
\end{document}
